\documentclass[11pt, oneside]{article} 
\usepackage{geometry}
\geometry{letterpaper} 
\usepackage{graphicx}
	
\usepackage{amssymb}
\usepackage{amsmath}
\usepackage{parskip}
\usepackage{color}
\usepackage{hyperref}

\graphicspath{{figures}{/Users/telliott/Github-Math/figures/}}

% \begin{center} \includegraphics [scale=0.4] {gauss3.png} \end{center}

\title{Rotation}
\date{}

\begin{document}
\maketitle
\Large

How to rotate a vector through an angle $\theta$ (counter-clockwise)?  Probably the simplest way is to memorize the matrix for rotation, and then test it by rotating unit vectors through $90^{\circ}$.  The matrix is

\[
\begin{bmatrix}  
\cos \theta & -\sin \theta \\
\sin \theta & \ \  \cos \theta
\end{bmatrix}
\]
Test it
\[
\begin{bmatrix}  
0 & -1 \\
1 & \ \ 0 
\end{bmatrix}
\begin{bmatrix}  1 \\ 0 \end{bmatrix}
=
\begin{bmatrix}  0 \\ 1 \end{bmatrix}
\ \ \ \ \ \ \ \ \ \ \ 
\begin{bmatrix}  
0 & -1 \\
1 & \ \ 0 
\end{bmatrix}
\begin{bmatrix}  0 \\ 1 \end{bmatrix}
=
\begin{bmatrix}  -1 \\ \ \ 0 \end{bmatrix}
\]
That looks good.  So
\[
\begin{bmatrix}  
\cos \theta & -\sin \theta \\
\sin \theta & \ \  \cos \theta 
\end{bmatrix}
\begin{bmatrix}  x \\ y \end{bmatrix}
=
\begin{bmatrix}  u \\ v \end{bmatrix}
\]
means that
\[ u = x \cos \theta - y \sin \theta \]
\[ v = x \sin \theta + y \cos \theta \]
Rotating the points counter-clockwise is equivalent to rotating the coordinate system clockwise.

\subsection*{rotation by $+ \theta$ then $- \theta$}

\[
\begin{bmatrix}  
\cos \theta & -\sin \theta \\
\sin \theta & \ \  \cos \theta
\end{bmatrix}
\begin{bmatrix}  
\ \ \cos \theta & \sin \theta \\
- \sin \theta &  \cos \theta
\end{bmatrix}
=
\begin{bmatrix}  
1 & 0 \\
0 & 1
\end{bmatrix}
\]

\subsection*{standard derivation}
Here is a second derivation with a sort of minimalist diagram of rotation
\begin{center} \includegraphics [scale=0.4] {min_rotation.png} \end{center}

We draw the horizontal $x$-axis and the rotated $u$-axis.  The angle between them is $t$.  We plot our point and then draw perpendiculars to both axes.  To finish the set-up, we draw perpendiculars from the point $(x,0)$ as shown.

Once the drawing is rendered, we are almost done.  Now we just work our way through.
\[ u = u_1 + u_2 = x \cos t + y \sin t \]

All the triangles in the diagram are similar, with small angle $t$.
\[ v = -v_2 + v_1 = -  x \sin t + y \cos t \]

I find this hard to remember.  You will know that you've done it right if you have both $x$ and $y$ as the \emph{hypotenuse of a right triangle}.  Start with the one for $x$ and then rotate the triangle the other direction about the point $(x,0)$.  Stretch to fit.

Here we are rotating the coordinate system counter-clockwise.  Hence, the minus sign is in the formula for the \emph{vertical} component, $v$.

\subsection*{Stewart}
I found a nice method in Stewart that depends on knowing the sum of angles formulas, which I urge you to memorize.
\begin{center} \includegraphics [scale=0.4] {min_rotation3.png} \end{center}

Here, we have the standard $xy$-coordinates in black.  The rotated $uv$-coordinate system is in red, with a rotation angle $\theta$ in the counter-clockwise direction.

The ray to the point $P$ has length $r$.  The coordinates of point $P$ in the $u,v$ system are naturally expressed in terms of $\phi$:
\[ u = r \cos \phi \]
\[  v = r \sin \phi \]
while $x$ and $y$ are naturally expressed in terms of the combined angle $\theta + \phi$.
\[ x = r \cos (\theta + \phi) \]
Now, use the sum formula for cosine: 
\[ x = r \cos \theta \cos \phi - r \sin \theta \sin \phi \]
But we have from above that
\[ u = r \cos \phi \]
\[ v = r \sin \phi \]
So
\[ x = u \cos \theta - v \sin \theta \]
It's as easy as that.

In the same way:
\[ y = r \sin (\theta + \phi) \]
\[ = r \sin \theta \cos \phi + r \cos \theta \sin \phi \]
\[ = u \sin \theta + v \cos \theta \]

The only trick here is that we have $x,y$ in terms of $u,v$.

\subsection*{Solve for u and v}
To convert these formulas to functions $u = f(x,y)$ and $v = f(x,y)$, there is a hard way and an easy way.  The easy way is to switch $x,y$ for $u,v$ and at the same time, substitute $-\theta$ for $\theta$.  What this amounts to is relabeling our diagram with $x,y$ being the rotated axes, and then rotating in the opposite direction (cw rather than ccw).
\[ x = u \cos \theta - v \sin \theta \]
switch
\[ u = x \cos -\theta - y \sin -\theta \]
\[ = x \cos \theta + y \sin \theta \]
(Recall that $\cos - x = \cos x$ and $\sin - x = - \sin x$).

For $y$
\[ y = u \sin \theta + v \cos \theta \]
switch
\[ v = x \sin -\theta + y \cos -\theta \]
\[ v = -x \sin \theta + y \cos \theta \]

\subsection*{examples}
Consider $y = x^2$ and rotate through an angle of $45^{\circ}$ cw.  The points are rotated cw and the coordinate system ccw, so the minus sign is in the expression for $x'$.

\[ x' = 1/\sqrt{2} (x - y) \]
\[ y' = 1/\sqrt{2} (x + y) \]
so
\[ y' = x'^2 \]
\[ 1/\sqrt{2} \cdot (x + y) = 1/2  \cdot (x - y)^2 \]
\[ x + y = 1/\sqrt{2}  \cdot (x - y)^2 \]
\begin{center} \includegraphics [scale=0.4] {rot_para.png} \end{center}

Here's another form --- both rotated and translated.
\begin{center} \includegraphics [scale=0.4] {quad4.png} \end{center}

Going back to our first result:
\[ x' = x \cos \theta - y \sin \theta \]
\[ y' = x \sin \theta + y \cos \theta \]

Let $x'y' = 1$, that is, the graph is a hyperbola.  We turn through an angle of $45^{\circ}$ cw.
\[ (x \cos \theta - y \sin \theta)(x \sin \theta + y \cos \theta) \]
\[ \frac{1}{2} \cdot (x - y)(x + y) \]
\[ 2  = (x^2 - y^2) \]
\begin{center} \includegraphics [scale=0.4] {hyper2.png} \end{center}
When this guy is rotated by $45$ ccw, it becomes the familiar graph of $xy=1$.
\begin{center} \includegraphics [scale=0.5] {hyper3.png} \end{center}

\end{document}
